\section*{\LARGE Adventure Background}\stepcounter{section}\phantomsection\addcontentsline{toc}{section}{Adventure Background}
{\entryfont Long before the modern Kingdom of Fife arose, dwarven settlements stretched beneath much of the land. Vast subterranean roads connected several great strongholds, allowing goods, knowledge, and travellers to move between distant parts of the dwarven realms without ever setting foot beneath the open sky. Although many smaller settlements existed throughout this network, three great centres became renowned for their individual crafts: Methven for alchemy, Aberdeen for smithwork and metallurgy, and Aberdour for constructs, weaponry, and experimental sources of energy.

Over time, however, relations between these settlements deteriorated. Their different traditions grew increasingly insular, and cooperation gave way to rivalry and distrust. Aberdour became particularly controversial among the other dwarven realms. Its engineers had begun constructing increasingly sophisticated mechanical servants, initially little more than simple labourers before gradually developing constructs capable of exploration, complex craftsmanship, and eventually warfare.

Seeking to push their inventions still further, the dwarves of Aberdour invited gnomish tinkers and engineers to live and work alongside them. Their aptitude for intricate mechanisms proved invaluable, and dwarven engineering soon began to incorporate gnomish ideas, symbols, and methods. To the other dwarven settlements, however, this cooperation represented an unacceptable breach of tradition. Outsiders had not merely been welcomed into Aberdour—they had been permitted to share in knowledge once guarded closely by the dwarves.

Whatever remained of the old unity was eventually shattered during the Great Eagle War between Angus and Fyfdonia. Tremors and devastation from the conflict collapsed great stretches of the subterranean network, while other passages were destroyed entirely. With its principal routes severed, the underground kingdom of Aberdour became completely isolated.

In the centuries that followed, even knowledge of its existence faded. Only scattered remnants of its technology survived close enough to the surface to be rediscovered. These vestiges would eventually contribute to the rise of modern Auchtertool and its extraordinary mechanical traditions, though their true origin was long forgotten.}

\subsection*{The Citadel of Aberdour}\stepcounter{subsection}\phantomsection\addcontentsline{toc}{subsection}{The Citadel of Aberdour}
{\entryfont The passage uncovered beneath Auchtertool's central plaza is not itself part of the lost city. Instead, it is one surviving branch of the ancient transportation network that once connected Aberdour with the other dwarven settlements beneath Fife.

Far beyond the newly opened entrance lies the Citadel of Aberdour, an immense subterranean complex that once stood at the centre of the kingdom's technological development. Within its workshops and laboratories, dwarven and gnomish engineers experimented with constructs, weapons, machinery, and forms of energy far beyond anything preserved in modern histories.

Most of the connecting passages have collapsed over the centuries, leaving the citadel effectively sealed away. The route discovered beneath Auchtertool is one of the rare surviving approaches.

Importantly, opening the passage beneath the plaza does not awaken the citadel or reactivate its defences. The newly exposed entrance is merely part of a dormant transportation route. Only once the expedition reaches the deeper reaches of Aberdour will the ancient mechanisms and guardians that remain within become relevant.}

\subsection*{The Energy Cores}\stepcounter{subsection}\phantomsection\addcontentsline{toc}{subsection}{The Energy Cores}
{\entryfont Much of Aberdour was once powered by enormous devices known only to modern observers as energy cores. Several remain embedded throughout the citadel, though most are inaccessible, inseparable from the machinery they power, or located far beyond the areas explored during this adventure.

The core discovered by the expedition is different.

It lies within what appears to have once been an engineering or assembly chamber and seems never to have been fully completed. Its outer structure remains unfinished, several components have yet to be properly integrated, and the device lacks whatever final refinements were required to make it fully operational. Nevertheless, its construction suggests a level of technological sophistication far beyond anything currently produced in Auchtertool.

Ancient guardian constructs still protect the chamber in which it rests. Whether they were assigned specifically to safeguard the unfinished core or simply to defend the workshop and its contents has long since been forgotten.

Recovering the device is possible, but completing or activating it is beyond the scope of this adventure. Whatever potential lies within the core will require considerably more study before it can be put to use.}

\subsection*{Ralathor-38B}\stepcounter{subsection}\phantomsection\addcontentsline{toc}{subsection}{Ralathor-38B}
{\entryfont Unknown to nearly everyone in this timeline, Ralathor-38B originates from another reality and has travelled across both time and the multiverse. The circumstances of his arrival and the means by which he accomplished this are not relevant to this adventure. To those who encounter him, he is known simply as the Mysterious Hermit.

Ralathor has come to this timeline for one purpose: to aid in the eventual defeat of Zargothrax. His knowledge of the forgotten dwarven tunnel network is central to that goal. Somewhere within those ancient passages lies a covert route into Dundee, one that Zargothrax himself does not know exists. Ralathor intends to reveal this route to Prince Angus McFife and the Questlords of Inverness when the proper moment arrives.}